\section{Teoria de Provas}

    \begin{itemize}[left=0.5cm, align=left, nosep]
        \item Um \textit{axioma} é uma fórmula bem-formada que representa uma proposição que é, geralmente, aceita como verdadeira. 
        \item Uma \textit{regra de inferência} é um mecanismo para obtenção de fôrmulas bem-formadas a partir de um conjunto de fórmulas bem-formadas.
        \item Sejam $\Gamma = \{ \gamma_1, \ldots, \gamma_n \} \subseteq \text{FBF}_{\mathcal{L}}, n \in \mathbb{N}, n \geq 0$ e $\varphi \in \text{FBF}_{\mathcal{L}}$.
        Regras de inferência são, normalmente, escritas da seguinte forma:
        
        \[
            \frac{
            \begin{array}{c}
            \gamma_1 \\
            \gamma_2 \\
            \vdots \\
            \gamma_n
            \end{array}
            }{\varphi}
        \]
        
        onde cada $\gamma_i \in \Gamma$ é chamada de \textit{premissa} e $\varphi$ é chamada de \textit{conclusão}. A aplicação da regra de inferência
        a $\gamma_1, \ldots, \gamma_n$ produz $\varphi$; nós também dizemos que $\varphi$ é \textit{derivada} de $\gamma_1, \ldots, \gamma_n$.    
    
    \end{itemize}

    \textbf{Exemplo 1:} 

    \begin{itemize}[left=0.5cm, nosep, label=$\hookrightarrow$]
        \item
    \end{itemize}

    \subsection{Propriedades Metateóricas} 
        
    $\{ p \to q, \neg q \} \models_{\mathcal{L}_p} \neg p$.
$\varphi \in \text{FBF}_{\mathcal{L}_P}$,
$\Gamma \models_{\mathcal{L}} \varphi$,
        \subsubsection*{}
$(\varphi \land \psi) \land \chi \Lpbimodels \varphi \land (\psi \land \chi)$
